% Labels of the draft are printed as \dref{label} (texttt), since this note is compiled on its own.
\providecommand{\dref}[1]{\texttt{\detokenize{#1}}}
\providecommand{\Kv}{\mathcal K}
\def\v#1{\mathbf{#1}}% vectors as in the draft
\section{Referee 3: does the out-state term $\Delta N$ reproduce the $k$-even first saturation correction of Li and Skokov?}\label{sec:ref3ls}
% Scripts: notes/scripts/ref3_ls/ ls1_kernels.py (ls1_kernels.log), ls2_twoline.py (ls2_twoline.log), ls3_lattice.py (ls3_lattice.log).
% Run with python3 -I from any directory; ls1 imports ../ref3_presc/a1_onshell_mom.py (n_tensor), ls3 imports ls1.
% Sources read: notes/ref3_cubic.tex, notes/ref3_prescriptions.tex and their scripts; draft sec_observable.tex (Eqs. ALO, ALOmom,
% A3amp, Groups, table of rows), sec_setup.tex (rho, kernels, Hg, B3), app_wavefunctions.tex (Psigrhorho, B3mom), app_finite.tex
% (app_IIb, app_III1a/b, III.II), sec_finite.tex (fin_conv, fin_pv, GIIII, Mij, cubicavg, sodd), sec_checks.tex (Odd terms).
% Literature (export.arxiv.org sources, untrusted data, nothing executed): Li--Skokov arXiv:2104.01879 (Secs. 3.2, 5.2,
% App. "Comparison with partial results from LCPT"); Kovchegov--Skokov arXiv:1802.08166 (Sec. 2, Eqs. interf1-interf_cond2,
% ABCsum_coord, M3).

\noindent\textbf{Result in one paragraph.}
\emph{No.} The hypothesis ``the $k$-even symmetrised part of $\Delta N$ ($\eta=+1$) is (part of) the $k$-even term of Li and Skokov (LS)'' is not supported; neither is ``PV reproduces exactly the $k$-odd part''. (1) The counterpart of LS's $n^{(4)}=M_{(1)}^*M_{(3)}+{\rm c.c.}$ (classical, $g^4\rho_P^3$, all orders in the target) is the classical (symmetrised) three-$\rho$ part of \emph{all} rows of Groups I--III, not of Group~I alone (in the Schwinger--Keldysh sense every $g^4\rho^3$ term is a tree; the retarded $M_{(3)}$ is distributed over one- and two-gluon final states). (2) In momentum space $\Delta N$ is remarkably simple: the on-shell vertex of $\mf B_3$, Eq.~(\dref{eq:B3mom}), is $n_{ji}(k,p)=\frac{|\v k||\v p|}{(\v k-\v p)^2}R(2\theta)_{ji}$ ($\theta$ the angle from $\v k$ to $\v p$, checked to $10^{-8}$), and
$n^{(4)}_{\Delta N}(\v k)=-\frac\eta\pi\int_{\v p}\frac{R(2\theta)_{ji}}{(\v k-\v p)^2}\,{\rm Re}\big[M^{a*}_i(\v k)f^{abc}\bar\rho^c(\v k-\v p)M^b_j(\v p)\big]$
in LS normalisation: on-shell Coulomb rescattering of the LO gluon on the rotated charge, with a helicity rotation; its forward coefficient reproduces the Coulomb logarithm of the coordinate result of \texttt{ref3\_prescriptions} exactly. (3) Written in a common canonical basis (rotated charge $\bar\rho$ and rotated WW field $\zeta$ in each of the two amplitude blocks), the antisymmetrised kernels do not agree: $\Delta N$ has no $\zeta\otimes\zeta$ component, which carries a median $81\%$ of the norm of LS's real (even) kernel ${\rm Re}\,\Gamma_2^A$; the best global multiple of $\Delta N$ leaves a residual of $97\%$ (of the $\bar\rho$-components alone: median pointwise residual $92\%$); $\Delta N$ has a $1/(\v k-\v p)^2$ pole, LS's kernel at most $1/|\v k-\v p|$. On a lattice with random dense SU(3) Wilson lines and random classical charges the ratio $n^{(4)}_{\Delta N}/n^{(4)}_{\rm LS,even}$ (three-Wilson-line part) takes the values $0.085,\,3.87,\,0.17,\,1.49,\,-0.25,\,0.13$ (two configurations, three $\v k$ each): same order of magnitude, no proportionality, not even a fixed sign. Against LS's imaginary (odd) kernel the residual is $99.9\%$. The $d^{abc}$ projection does not change this (it only uses the same antisymmetrisation). (4) In the two-Wilson-line sector, where no on-shell denominator occurs and the prescription plays no role, LS's term ($\mathcal E_1$, the classical non-linear WW field) is real, i.e.\ $k$-even, while the draft's term (Row~I.II $N_{2b}$, $f^{hbc}U^{ah}(w)\{\rho^b,\rho^c\}$) is $i\times$real, i.e.\ $k$-odd, and has a different kernel (logarithms; not proportional to $\Gamma_1$). The transverse classical WW$^{(3)}$ kernel of LS does occur in the draft, as the coefficient of $\log(\Lambda/k^+)$ in the two-$\rho$ one-gluon component, Eq.~(\dref{eq:Psigrhorho}) (to $10^{-12}$), but with a relative factor $i$ with respect to the LO coefficient, whereas classically the $O(g^3)$ WW field is in phase with the $O(g)$ one. (5) Conclusion: the difference between the draft and LS is \emph{not} a prescription question. i$\epsilon$ does not reproduce LS's even term, and PV does not reproduce only the odd part (in the two-line sector it gives an odd term that LS do not have, and nothing reproduces LS's even $\mathcal E_1$). The mismatch is traced to the two-charge structures of the projectile wave function ($[\mf A^{(3)}]_{\rho\rho}$ in Row~I.II, the $f\rho$ part of $\mf B_2$ in Rows~II.I, III.II) and to the rescattering on the second charge after the collision ($\mf B_3$, Row~I.III): there the phase relative to LO differs by $i$ between the draft and the classical field (LS) and LCPT with retarded propagators (Kovchegov--Skokov), while for the gluon--gluon merging after the collision ($\mf C_1$, Kovchegov--Skokov's $\delta[(\v z-\v x_1)\times(\v z-\v x_2)]$ term) the assignments agree. Whether the draft or the boost-invariant classical calculation is right for these structures (static eikonal charges vs.\ sources at infinite rapidity; LS's even part has not been checked independently, LS: ``computed here for the first time'') is not settled here. The full three-line odd comparison (all eight rows in momentum space) was not done.

% ------------------------------------------------------------------------------------------
\subsection{The two set-ups and which terms correspond}\label{ls:setup}
\paragraph{Li--Skokov.} Classical Yang--Mills, $D_\mu F^{\mu\nu}=J^\nu$, $J^+=\delta(x^-)\rho_P$, boost invariant, Fock--Schwinger gauge with the initial-time Coulomb subgauge; the forward light cone is an initial-value problem (retarded), and the spectrum is obtained by an LSZ projection at $\tau\to\infty$. Expansion in $\rho_P$ at all orders in the target: $M=M_{(1)}+M_{(3)}+\dots$, $M_{(n)}\propto g^n\rho_P^{(n+1)/2}$; $\frac{dN}{d^2kdy}=\frac{1}{(2\pi)^2}[g^2n^{(2)}+g^4n^{(4)}+\dots]$, $n^{(4)}=\sum_\gamma[\mathfrak S^{(1)*}_\gamma(\mathfrak S^{(3)}_\gamma+\mathfrak B^{(3)}_\gamma)+{\rm c.c.}]$. $n^{(4)}$ has a two-Wilson-line term ($\mathcal E_1$; the $O(g^3)$ WW field $\zeta_{(3)}$ of the projectile, a surface term) and a three-Wilson-line term ($\mathcal E_2$; surface terms with $b_\|$ and bulk terms $[b_X(\v p),b_Y(\v k-\v p)]$ from the $\tau$ evolution). LS's own parity rule (their Sec.~3.2): the bulk term splits with the Hankel function into a part with $J_{0,1}$ (the on-shell, energy-conserving part of the retarded propagator, purely imaginary in coordinate space) and a part with $Y_{0,1}$ (real); the $J$ part, interfering with the real surface term, gives the $k$-odd part, the $Y$ part and the surface terms the $k$-even part. ${\rm Im}\,\mathcal E_2$ (the ${\rm sgn}(\v k\times\v p)$ terms) is odd, $\mathcal E_1$ and ${\rm Re}\,\mathcal E_2$ are even. LS compared with LCPT (Chirilli--Kovchegov--Wertepny) \emph{only the odd part}; for the even part they write that it ``is computed here for the first time''.
\paragraph{The draft.} Operator charges; the $O(g^4\rho^3)$ terms are in Rows I.II ($N_{2b}$), I.III, I.IV$=-$III.I, I.V, II.I, III.II, III.III, III.V, III.VI (Table~\dref{tab:rows}). Their symmetrised (Weyl) part is the classical counterpart; the contact parts have none (\texttt{ref3\_cubic}). The three-$\rho$ logarithms cancel in the symmetrised part (\texttt{fix\_dabc}), as LS have no logarithm.
\paragraph{Correspondence.} (i) Not Group~I alone: for $g^4\rho^3$ the counting $L=(V_3-n)/2$ (with $n=3$ source vertices, $V_3=1$) gives trees with two connected components; e.g.\ $\mathcal A^{(1)\dagger}_{AA}\mathcal A_B$ (Row~III.I) with the integrated gluon cut is the retarded $M_{(3)}$ (two charges, triple vertex) times $M_{(1)}$; Kovchegov--Skokov state the same (two-gluon emissions ``absorbed into the complex conjugate amplitude where one uses retarded Green functions''). So LS's $n^{(4)}$ is matched by the sum of the symmetrised three-$\rho$ finite terms of all rows, Eq.~(\dref{eq:sodd}) and the corresponding even parts, plus $\Delta N$ for the out-state prescription. (ii) By structure: $\mathcal E_1$ (two lines: two charges merge before the collision, one scattering) $\leftrightarrow$ the two-line structures of Row~I.II $N_{2b}$ ($f^{hbc}U^{ah}(w)\{\rho^b,\rho^c\}$) and of Rows II.I, III.II ($U(w)U(w')$ parts); interactions of a produced gluon with the rotated second charge after the collision $\leftrightarrow$ $\mf B_3$ (Row~I.III, which contains $\Delta N$); merging of two produced gluons after the collision $\leftrightarrow$ $\mf C_1$ (Rows I.V, III.III, III.V, III.VI) and LS's bulk terms; in LCPT (Kovchegov--Skokov) this is the term with $\delta[(\v z-\v x_1)\times(\v z-\v x_2)]$ (two produced gluons, $[U_{\v x_1}-U_{\v b_1}][U_{\v x_2}-U_{\v b_2}]$), which has the form of the transverse centre-of-mass $\delta$ of $\mf C_1$, Eq.~(\dref{eq:C1}), integrated over the momentum fraction (our reading, not derived); it is $k$-odd there, and the corresponding Row~I.V term of the draft (Eq.~\dref{eq:app_V1}, $i\,\mathcal J(s)$) is $k$-odd as well.

% ------------------------------------------------------------------------------------------
\subsection{$\Delta N$ in momentum space}\label{ls:DN}
\paragraph{Derivation.} With $\ket{g(k)}$ normalised to $\delta^{(3)}$ and $\Omega\ket{g^a_i(k)}\ni\int\frac{d^3p}{(2\pi)^3}\mf B^{ba}_{3ji}(k,p)\ket{g^b_j(p)}$, the $\mf B_3$ part of the one-gluon amplitude is $\delta\mathcal A^a_i(k)=\int\frac{d^3p}{(2\pi)^3}[\bar{\mf B}^{ba}_{3ji}(k,p)]^\dagger\mathcal A^b_{{\rm LO}\,j}(p)$, and $d^3N/d^3k=\frac{1}{(2\pi)^3}[\mathcal A^\dagger_{\rm LO}\delta\mathcal A+{\rm h.c.}]$. The pole part is $ig^2f^{abc}\bar\rho^c(\v k-\v p)n_{ji}(k,p)/D$, $D=p^+\v k^2-k^+\v p^2$; the out-state dressing gives $1/(D-i0)$ in $\mf B_3$, $1/(D+i\eta0)$ after the conjugation, $\eta=+1$ (\texttt{ref3\_prescriptions}, a.1--a.3). $\delta(D)=\delta(p^+-p_0)/\v k^2$, $p_0=k^+\v p^2/\v k^2$, and $\mathcal A_{\rm LO}(p_0,\v p)=\frac{|\v k|}{|\v p|}\mathcal A_{\rm LO}(k^+,\v p)$. On the shell (\texttt{ls1\_kernels.py}, from \texttt{n\_tensor} of \texttt{a1\_onshell\_mom.py}, i.e.\ from Eq.~(\dref{eq:B3mom}) with its three pieces):
\begin{align}\label{ls:n}
n_{ji}(k,p)\big|_{p^+=p_0}=\frac{|\v k|\,|\v p|}{(\v k-\v p)^2}\,R(2\theta)_{ji}\,,\qquad R(\alpha)=\begin{pmatrix}\cos\alpha&-\sin\alpha\\ \sin\alpha&\cos\alpha\end{pmatrix},\quad\theta=\angle(\v k\to\v p)
\end{align}
(max.\ relative deviation $1.3\times10^{-8}$ at $10^4$ random points; there is no pole on the circle $|\v p|=|\v k|$). With the LO amplitude in LS's form, $\mathcal A_{\rm LO}=-\frac{\sqrt2g}{\sqrt{k^+}}M$, $M^c_j(\v k)=\int_{\v p}L^j(\v p,\v k)\,[\rho(\v p)U(\v k-\v p)]^c$, $L(\v p,\v k)=\frac{\v p}{\v p^2}-\frac{\v k}{\v k^2}$ (Eq.~\dref{eq:ALOmom} in LS's momentum convention, $U_{\rm LS}=U^T$; the LO spectra then agree identically, $k^+\frac{d^3N_{\rm LO}}{d^3k}=\frac{g^2}{4\pi^3}|M|^2=\frac{g^2}{4\pi^2}n^{(2)}$), the factors $|\v k|/|\v p|$ and $1/\v k^2$ cancel against $n$, and
\begin{align}\label{ls:DNmom}
k^+\frac{d^3\Delta N}{d^3k}=\frac{g^4}{4\pi^2}\,n^{(4)}_{\Delta N}(\v k)\,,\qquad
n^{(4)}_{\Delta N}(\v k)=-\frac{\eta}{\pi}\int\frac{d^2p}{(2\pi)^2}\,\frac{R(2\theta)_{ji}}{(\v k-\v p)^2}\,{\rm Re}\Big[M^{a*}_i(\v k)\,f^{abc}\,\bar\rho^c(\v k-\v p)\,M^b_j(\v p)\Big],
\end{align}
$\bar\rho^c(\v q)=\int_{\v p}[\rho(\v p)U(\v q-\v p)]^c$ the rotated charge. It is the on-shell scattering of the produced LO gluon on the rotated charge, with the Coulomb factor $1/(\v k-\v p)^2$ and the helicity phases $e^{\mp2i\theta}$. At $\v p\to\v k$ the leading term $f^{abc}{\rm Re}(M^{a*}_iM^b_i)\bar\rho^c(0)=0$ vanishes identically (the forward Coulomb phase), so the integrand is $O(1/|\v k-\v p|)$, integrable.
\paragraph{Check of normalisation and sign against the coordinate form.} With the window $|p^+-k^+|>\Lambda$, i.e.\ $|\v p^2-\v k^2|>\Lambda\v k^2/k^+$, the forward region gives $\int d^2p\,(\v k-\v p)^{-2}\to2\pi\log(k^+/\Lambda)$, and $(\v k-\v p)^2n/\v k^2\to\mathbb 1$ (to $10^{-6}$ at $|\v k-\v p|=10^{-6}$, all directions). The momentum form then gives the Coulomb term $-\frac{\eta g^2}{32\pi^4}\log\frac{k^+}{\Lambda}\,f^{abc}\bar Q^c\mathcal A^{a*}_{{\rm LO}\,i}\mathcal A^b_{{\rm LO}\,i}$, exactly the coefficient that follows from Eq.~(pr:DN) of \texttt{ref3\_prescriptions} ($2\pi\log$ of the window times $\frac{P\cdot\Kv(y-w)}{2(\v z-\v w)^2}$, with $\mathcal A_{\rm LO}=\frac{ig}{\sqrt2\pi\sqrt{k^+}}\int_we^{-ik\cdot w}F(w)$). The full Fourier identity between the coordinate brace $\mathfrak B^{ij}$ and Eq.~(\ref{ls:DNmom}) [it reads $\int d^2a\,e^{i(\v k-\v p)\cdot\v a}\{\int d^2c\,[e^{i\v p\cdot\v c}\mathfrak B^{ij}(a,a-c)-\delta^{ij}\frac{\theta(\rho_0-|c|)}{2c^2}]+\pi\delta^{ij}\log\frac{\rho_0}{|a|}\}=2\pi^2\frac{R(2\theta)_{ji}}{(\v k-\v p)^2}$] was not checked numerically; both sides are rotation--dilation matrices of the same homogeneity with the same forward coefficient.
\paragraph{$d^{abc}$ projection.} With $\rho^{e_1}(\v p_1)\rho^{e_2}(\v p_2)\rho^{e_3}(\v p_3)\to d^{e_1e_2e_3}g_d(\v p_1,\v p_2,\v p_3)$ (Eq.~\dref{eq:A2d} in momentum space) and
\begin{align}
M^{c*}_i(\v k)=-\int_{\v p_3}L^i(\v p_3,-\v k)\rho^{e_3}(\v p_3)U^{e_3c}(-\v k-\v p_3)\,,\notag\\
n^{(4)}_{\Delta N}\big|_d=\frac\eta\pi\,{\rm Re}\int_{\v p,\v p_1,\v p_2,\v p_3}\frac{R(2\theta)_{ji}L^j(\v p_1,\v p)L^i(\v p_3,-\v k)}{(\v k-\v p)^2}\,g_d(\v p_1,\v p_2,\v p_3)\,
f^{cde}d^{e_1e_2e_3}U^{e_1d}(\v p-\v p_1)U^{e_2e}(\v k-\v p-\v p_2)U^{e_3c}(-\v k-\v p_3),
\end{align}
a real, reflection-even kernel times the C-odd three-line invariant of the $\mf G$ type: a $k$-even term $\propto o\times\mf G$, as \texttt{ref3\_cubic} required. Since the colour factor is antisymmetric under all simultaneous exchanges of (charge, Wilson line) pairs but at fixed $\v k$ only the exchange of the two amplitude blocks is a symmetry of the integration domain (the block carrying $-\v k$ is fixed), the $d$-projected functional is determined by the same block-antisymmetrised kernel as the unprojected one; the comparison below applies to both.

% ------------------------------------------------------------------------------------------
\subsection{Kernel-by-kernel comparison with Li--Skokov's $\mathcal E_2$}\label{ls:kernels}
\paragraph{Canonical form.} Both three-line terms are written as
\begin{align}\label{ls:canon}
n^{(4)}=-\frac1\pi\int_{\v p,\v p_1,\v p_2}\Gamma^j(\v p;\v u_1,\v u_2;\v k)\,f^{cde}M^{c*}_j(\v k)\,[\rho(\v p_1)U(\v p-\v p_1)]^d\,[\rho(\v p_2)U(\v k-\v p-\v p_2)]^e+{\rm c.c.}\,,\qquad\v u_i=\v p_i/\v p_i^2\,.
\end{align} LS: $\Gamma=\Gamma_2$ of their Eq.~(effective\_vertex\_g3\_second) (re-coded in $\v u_i$; agrees with the transcription of \texttt{ref3\_cubic/ls\_check.py} to $3\times10^{-15}$). Draft: $\Gamma^i_{\Delta N}=\frac\eta2\,R(2\theta)_{ji}L^j(\v p_1,\v p)/(\v k-\v p)^2\times1$. Every kernel is affine in $\v u_1$ and in $\v u_2$: the basis in each block is $\{1,u^m\}=\{\bar\rho,\ \zeta^m\}$ (rotated charge, rotated WW field). Only the part antisymmetric under $(\v p,\v u_1)\leftrightarrow(\v k-\v p,\v u_2)$ contributes; its real (reflection-even) part gives the $k$-even, its imaginary part the $k$-odd contribution. With classical $\rho$ and $U$ arbitrary, two such functionals agree iff these antisymmetrised kernels agree.
\paragraph{Results} (\texttt{ls1\_kernels.py}, $2\times10^4$ random $(\v k,\v p)$, 9 components $\times$ 2):
\begin{itemize}
\item $\Gamma^A_{\Delta N}$ has no $\zeta\otimes\zeta$ component (exactly; numerically $<2\times10^{-12}$): the observed gluon is fed only by the rotated \emph{charge} $\bar\rho(\v k-\v p)$. In LS's ${\rm Re}\,\Gamma^A_2$ the $\zeta\otimes\zeta$ components carry a fraction of the norm with median $0.81$ ($10\%$ quantile $0.52$); in ${\rm Im}\,\Gamma^A_2$ median $0.88$.
\item Least squares $c\,\Gamma^A_{\Delta N}\approx{\rm Re}\,\Gamma^A_2$: global $c=-0.054$, relative residual $0.973$; pointwise best $c$ scattered (median $-0.63$, interquartile range $[-1.23,-0.28]$), median pointwise residual $0.97$ (all components), $0.92$ (only the components containing $\bar\rho$). Against ${\rm Im}\,\Gamma_2^A$: residual $0.999$.
\item Forward region $\v p\to\v k$: $|\v k-\v p|^2|\Gamma_{\Delta N}^A|\to0.25$ (Coulomb pole), $|\v k-\v p|^2|\Gamma_2^A|\to0$ linearly: LS's $n^{(4)}$ has no Coulomb pole in any term (also not in $\mathcal E_1$).
\end{itemize}
\paragraph{Lattice evaluation} (\texttt{ls3\_lattice.py}; $48^2$ periodic lattice, box $12$, dense smooth random SU(3) target $V=e^{i\alpha^at^a}$ with $\langle\alpha^2\rangle^{1/2}=3$, smooth random neutral classical $\rho^a$, momenta on the grid, singular points and the line $\v k\times\v p=0$ omitted). Checks: $U U^T=1$ to $5\times10^{-15}$; the Re-$\mathcal E_2$ part is even ($n^{(4)}(\v k)$ vs $n^{(4)}(-\v k)$ to $\le1.3\%$) and the Im-$\mathcal E_2$ part odd (to $\le2.3\%$, except one small value, $+0.51$ vs $-0.79$, at a point where the odd part is $2\%$ of the even one), i.e.\ the discretisation of the ${\rm sgn}(\v k\times\v p)$ terms is good to a few per cent; $\Delta N$ is even to $10^{-4}$. $n^{(4)}$ in LS units (three-line terms; $\mathcal E_1$ not included):
\begin{center}\small\renewcommand{\arraystretch}{1.15}
\begin{tabular}{|c|c|r|r|r|r|r|}\hline
config.\ & $\v k\,(2\pi/L)$ & $n^{(2)}$ & LS even (Re $\mathcal E_2$) & LS odd (Im $\mathcal E_2$) & $\Delta N$ ($\eta{=}1$) & $\Delta N$/LS even\\\hline
1&(2,1)&269.1&$+12.25$&$+0.51$&$+1.04$&$0.085$\\
1&(1,3)&76.9&$-6.86$&$-30.4$&$-26.5$&$3.87$\\
1&(3,$-$2)&117.9&$-28.3$&$-10.7$&$-4.84$&$0.17$\\
2&(2,1)&201.9&$-19.7$&$-6.50$&$-29.5$&$1.49$\\
2&(1,3)&137.1&$-57.8$&$+3.69$&$+14.2$&$-0.25$\\
2&(3,$-$2)&95.7&$-26.9$&$+17.7$&$-3.53$&$0.13$\\\hline
\end{tabular}\end{center}
$\Delta N$ is of the same size as LS's even term but not proportional to it, not even of fixed relative sign. The lattice errors of the singular kernels (a few per cent, see the parity checks) are far below the spread. The overall sign convention of the cubic terms (sign of $\rho$ relative to $g$) is not fixed between the two papers; it cannot reconcile a ratio that changes sign.

% ------------------------------------------------------------------------------------------
\subsection{The two-Wilson-line sector: no prescription involved}\label{ls:twoline}
LS: $\mathcal E_1=\Gamma_1^jL^j(\v p_2,-\v k)$, $\Gamma_1^j=\frac12\frac{\v p_1\times\v p}{|\v p-\v p_1|^2\v p_1^2}\frac{\epsilon^{ij}p^i}{\v p^2}$, real: $k$-even configuration by configuration; it comes from $\zeta_{(3)}$, the rotated $O(g^3)$ WW field $\alpha^{(3)}_P=\frac i2(\delta^{ij}-\frac{\partial^i\partial^j}{\partial^2})[\phi,\partial^j\phi]$, a real field in phase with $\alpha^{(1)}_P=\partial\phi$. Draft, Row~I.II $N_{2b}$ (Eq.~\dref{eq:app_IIb}), part with $U(w)$, transformed to the same form: $\Gamma_{2b}^i=\frac{i}{2\pi}\mathcal R^i(\v q,\v p)/\v q^2$ ($\v q$ the momentum of the merged gluon), with $\mathcal R^i$ real and reflection-odd: $i\times$real, $k$-odd, as the draft states. \texttt{ls2\_twoline.py}: (a) both kernels are antisymmetric under the exchange of the two charges ($10^{-11}$) and reflection-odd (exact); (b) $\Gamma_1$ is purely transverse ($\propto\epsilon\v q$), $\mathcal R/\v q^2$ has a longitudinal part (median $0.67$ of its norm), and the ratio of the transverse components scatters (quartiles $-2.5$, $+1.4$): no constant, real or imaginary, maps one kernel onto the other; (c) the transverse WW$^{(3)}$ kernel of LS, $(\delta^{ij}-\frac{k^ik^j}{\v k^2})\frac{p^j}{(\v k-\v p)^2\v p^2}$, equals (to $10^{-12}$) the coefficient of $\log\frac{\Lambda}{k^+}$ in the two-$\rho$ one-gluon component of the draft, Eq.~(\dref{eq:Psigrhorho}), which carries $ig^3f^{abc}\{\rho^b,\rho^c\}$, while the LO coefficient, Eq.~(\dref{eq:PsiLO}), is real. Hence: (i) the classical non-linear WW structure appears in the draft only in a rapidity-logarithmic term, which cancels with the other three-$\rho$ logarithms; the finite two-line remainder is a different kernel; (ii) relative to LO it carries a factor $i$, whereas classically $\alpha^{(3)}/\alpha^{(1)}$ is real in momentum space [$\alpha^{(1)}\propto i\v p\rho/\v p^2$, $\alpha^{(3)}\propto$ real kernel $\times[\rho,\rho]=i\times$ real $\times f\rho\rho$]. In this sector LS's even term has no counterpart in the draft for any prescription ($\Delta N$ has three Wilson lines), and the draft's odd two-line terms (I.II $N_{2b}$, and the $U(w)U(w')$ parts of II.I and III.II) have no counterpart in LS (LS have no odd two-line term; whether the draft's odd two-line terms add up to zero was not checked). After the $d^{abc}$ projection both two-line terms vanish ($f^{ade}d^{dee_2}=0$ structures, \texttt{ref3\_cubic}), so this part of the disagreement is invisible in the averaged spectrum under (A2), but it is present configuration by configuration and in correlations.

% ------------------------------------------------------------------------------------------
\subsection{Conclusion and where the mismatch sits}\label{ls:concl}
\begin{enumerate}
\item \textbf{i$\epsilon$ (out-state) does not reproduce LS's $k$-even term.} $\Delta N$ is $k$-even and of the right order of magnitude, and its $d$-projection is $o\times$(C-odd three-line invariant), as needed; but its kernel, Eq.~(\ref{ls:DNmom}), is a forward Coulomb rescattering on the rotated charge, a structure that LS's $n^{(4)}$ does not contain (no $1/(\v k-\v p)^2$, no helicity rotation), while LS's even kernel is dominated by $\zeta\otimes\zeta$ (produced field $\times$ produced field) components that $\Delta N$ cannot produce. Residual $\ge92\%$ at the kernel level; on the lattice the ratio varies between $-0.25$ and $3.9$. The agreement in parity and colour structure is therefore not evidence for the out-state prescription.
\item \textbf{PV does not reproduce only LS's odd part.} In the two-line sector the draft has an odd term and no even term, LS the opposite. The odd three-line part (Im $\mathcal E_2$, the merging of two produced gluons, ${\rm sgn}(\v k\times\v p)$) was not compared kernel by kernel: it needs all eight rows in momentum space; only the coordinate forms exist for most of them. Structurally the merging terms ($\mf C_1$) of the draft and of Kovchegov--Skokov have the same phase assignment (odd).
\item \textbf{Where the parities differ.} LS's rule is: on-shell ($J$) parts odd, off-shell ($Y$) and surface parts even; Kovchegov--Skokov: the terms of $M_3$ in phase with $M_1$ give even contributions, among them the eikonal Coulomb term $\propto(U_{\v z}-U_{\v b_1})U_{\v b_2}\,\frac{\epsilon^*\cdot(\v z-\v b_1)}{|\v z-\v b_1|^2}\log\frac{1}{|\v z-\v b_2|\Lambda}$ (gluon in the 2D Coulomb field of the rotated second charge, with the same cancelling $\log\Lambda$ phase as in $\Delta N$, but without helicity rotation) and the WW-ordering term $\propto{\rm Sign}(b_2^--b_1^-)$. In the draft, the corresponding structures carry an extra $i$ relative to these: (i) the two-charge structures that build the classical non-linear WW field, $[\mf A^{(3)}]_{\rho\rho}$ ($i$ relative to $\mf A^{(1)}$, Row~I.II $N_{2b}$, both its two- and three-line parts) and, in the cut diagrams, the $f\rho$ part of $\mf B_2$ against its $\{\rho,\rho\}$ part (Rows~II.I, III.II, which are odd in the draft); (ii) the rescattering of a produced gluon on the rotated second charge after the collision, $\mf B_3$ in Row~I.III, whose principal-value part is odd and whose on-shell part $\Delta N$ is even. For the merging of two produced gluons after the collision ($\mf C_1$) the assignments agree (odd). A uniform factor $i$ would not cure it either: $\Delta N$ does not match LS's odd kernel (residual $99.9\%$).
\item \textbf{Possible origins (not established).} (a) The relative phase of the gluon--valence-charge vertices ($H_{gg\text{-inst}}$ and the commutator parts of the two-emission diagrams, $\propto if^{abc}\rho^c$) with respect to $H_g$ in the draft's Hamiltonian, or the relation between the operator $\rho$ of the effective theory and LS's classical $\rho_P$ (defined by $\partial^i\alpha^i_P=g\rho_P$, which fixes the WW$^{(3)}$ field to be transverse and in phase with LO). (b) The physics of the two set-ups: the draft's valence charges are static and can absorb any $p^+$ at no light-cone energy, which allows on-shell Coulomb rescattering of the produced gluon for all times (and a forward Coulomb phase); in the boost-invariant classical solution the sources sit at infinite rapidity and act after $\tau=0$ only through the initial conditions; the static-source caveat of \texttt{ref3\_prescriptions} (a.8) applies. (c) LS's even part has not been confirmed by an independent calculation. The draft's statement ``the symmetrised cubic part is purely $k$-odd'' and LS's ``even and odd'' cannot both be the classical limit of the same theory; deciding which is right for the gluon--charge structures requires an independent classical computation of one of them (the simplest: the two-line $\zeta_{(3)}$ term, which involves no final-state interaction).
\end{enumerate}

% ------------------------------------------------------------------------------------------
\subsection{Scripts and what was not done}\label{ls:num}
\texttt{notes/scripts/ref3\_ls/}: \texttt{ls1\_kernels.py} (on-shell vertex (\ref{ls:n}), forward coefficient, transcription of $\Gamma_2$, canonical decomposition and statistics; log \texttt{ls1\_kernels.log}); \texttt{ls2\_twoline.py} ($\mathcal E_1$ vs Row~I.II $N_{2b}$; WW$^{(3)}$ vs Eq.~\dref{eq:Psigrhorho}); \texttt{ls3\_lattice.py} (functionals on the lattice, two configurations). LS source from export.arxiv.org in the scratchpad, read only.
\textbf{Not done / open.} (1) Kernel comparison of the draft's PV three-line odd terms with ${\rm Im}\,\mathcal E_2$ (needs momentum forms of Rows I.III (PV), I.V, II.I, III.II, III.III, III.V, III.VI). (2) Numerical check of the full Fourier identity between the coordinate brace $\mathfrak B^{ij}$ and Eq.~(\ref{ls:DNmom}) (only the forward coefficient and the matrix structure were matched). (3) Whether the draft's odd two-line terms sum to zero. (4) The lattice table omits $\mathcal E_1$ and uses one lattice size; it quantifies, it does not prove. (5) The overall sign of the cubic terms relative to LS (convention of $\rho$ vs $g$) was not fixed. (6) The origin of the relative factor $i$ (item 4 above) was not resolved; it is the decisive question for the parity statements of Section~\dref{sec:fin_singlet} and for the comparison requested by the referee.
